\documentclass[12pt,a4paper]{article}
\usepackage{fontspec}
\setmainfont[Path=fonts/,BoldFont=DejaVuSerif-Bold.ttf,ItalicFont=DejaVuSerif-Italic.ttf,BoldItalicFont=DejaVuSerif-BoldItalic.ttf]{DejaVuSerif.ttf}
\setmonofont[Path=fonts/]{DejaVuSansMono.ttf}
\usepackage[russian]{babel}
\usepackage{amsmath,amssymb,graphicx,booktabs,array,float}
\usepackage[left=25mm,right=20mm,top=22mm,bottom=24mm]{geometry}
\usepackage[unicode,hidelinks]{hyperref}
\usepackage{caption}
\captionsetup{font=small,labelsep=period}
\setlength{\parindent}{1.25cm}
\setlength{\parskip}{4pt}
\setlength{\emergencystretch}{3em}
\renewcommand{\arraystretch}{1.15}
\newcommand{\norm}[1]{\left\lVert #1\right\rVert}
\newcommand{\pic}[3]{\begin{figure}[H]\centering\includegraphics[width=#2\textwidth]{images/#1}\caption{#3}\end{figure}}
\begin{document}
\begin{titlepage}
\begin{center}
Московский государственный университет\\имени М.\,В.~Ломоносова\\[6mm]
Факультет вычислительной математики и кибернетики
\vfill
{\large Отчёт по практическому заданию\\по курсу\\[3mm]
«Суперкомпьютерное моделирование и технологии»}\\[15mm]
{\Large\bfseries Задача Дирихле для уравнения Пуассона\\[3mm]в криволинейной области}\\[8mm]
{\large Задача 1. Вариант 15\\[4mm]Этап 1. OpenMP-реализация}
\vfill
\begin{flushright}
Выполнил студент группы 609\\
Кабедев Вадим Дмитриевич
\end{flushright}
\vfill
Москва, 2026
\end{center}
\end{titlepage}
\setcounter{page}{2}

\section{Математическая постановка}
Требуется найти приближённое решение задачи
\[
-\Delta u=1\quad\text{в }D,\qquad
u=\varphi(x,y)=\sin(x+y)\quad\text{на }\partial D,
\]
где область варианта 15 имеет вид
\[
D=\{(x,y):x^2+(y-2)^2<9\}
\cap\{(x,y):x^2+(y+2)^2>9\}.
\]
Область находится внутри круга радиуса 3 с центром $(0,2)$ и вне круга того же радиуса с центром $(0,-2)$. Для метода фиктивных областей выбран прямоугольник
$\Pi=[-3,3]\times[0,5]$.
Продолжение граничного значения на прямоугольник:
$v_0(x,y)=\sin(x+y)$.

\section{Метод фиктивных областей и разностная схема}
В прямоугольнике рассматривается задача
\[
-\Delta v+Q(x,y)(v-v_0)=F(x,y),\qquad v|_{\partial\Pi}=v_0,
\]
где $Q=0$ в $D$, $Q=\varepsilon^{-2}$ в фиктивной области,
$F=1$ в $D$ и $F=0$ вне $D$.
Используется равномерная сетка
\[
x_i=-3+ih_1,\quad y_j=jh_2,\qquad
h_1=\frac6M,\quad h_2=\frac5N,\quad
\varepsilon=h^2,\quad h=\max(h_1,h_2).
\]
Числа $M,N$ обозначают интервалы. Полная сетка имеет
$(M+1)(N+1)$ узлов, число неизвестных равно $(M-1)(N-1)$.

Пусть $S_{ij}$ --- площадь пересечения области $D$ с ячейкой вокруг узла $(x_i,y_j)$, а
$\theta_{ij}=S_{ij}/(h_1h_2)$ --- доля площади. Тогда
\[
Q_{ij}=\frac{1-\theta_{ij}}{\varepsilon^2},
\qquad F_{ij}=\theta_{ij}.
\]
Доля ограничивается диапазоном $[0,1]$ для устранения выхода за него из-за округления. Целиком внутренние и внешние ячейки распознаются по расстояниям до центров окружностей.

Для граничных ячеек площадь вычисляется аналитическим интегрированием вертикальных сечений. При $s(x)=\sqrt{9-x^2}$ область описывается неравенствами
\[
|s(x)-2|<y<2+s(x),\qquad -3<x<3.
\]
Интервал интегрирования разбивается по пересечениям окружностей с рёбрами ячейки. На каждом участке длина сечения имеет вид $a+b\sqrt{9-x^2}$; используется первообразная
\[
\int\sqrt{9-x^2}\,dx
=\frac12\left(x\sqrt{9-x^2}+9\arcsin\frac x3\right).
\]

Разностный оператор задаётся пятиточечным шаблоном:
\begin{align*}
(Aw)_{ij}
&=d_{ij}w_{ij}
-\frac{w_{i-1,j}+w_{i+1,j}}{h_1^2}
-\frac{w_{i,j-1}+w_{i,j+1}}{h_2^2},\\
d_{ij}&=\frac2{h_1^2}+\frac2{h_2^2}+Q_{ij},\qquad
B_{ij}=F_{ij}+Q_{ij}\sin(x_i+y_j).
\end{align*}
Известные значения на внешней границе дополнительно переносятся в правую часть с коэффициентами $h_1^{-2}$ и $h_2^{-2}$. Векторы, к которым применяется оператор, имеют нулевые внешние граничные элементы. Матрица явно не хранится.

Оператор симметричен и положительно определён: его квадратичная форма представляет сумму квадратов разностных градиентов и неотрицательного слагаемого с $Q$ при однородных внешних условиях. Поэтому применим метод сопряжённых градиентов.

\section{Итерационный алгоритм}
Используются скалярное произведение и норма
\[
(a,b)=h_1h_2\sum_{i=1}^{M-1}\sum_{j=1}^{N-1}a_{ij}b_{ij},
\qquad \norm a_E=\sqrt{(a,a)}.
\]
Предобуславливатель $D_{\rm diag}$ равен диагонали матрицы $A$.
Начальное приближение $w^{(0)}=0$,
$r^{(0)}=B-Aw^{(0)}$, $z^{(0)}=D_{\rm diag}^{-1}r^{(0)}$, $p^{(1)}=z^{(0)}$.

На каждой итерации выполняются операции
\begin{align*}
\alpha_k&=\frac{(r^{(k-1)},z^{(k-1)})}{(Ap^{(k)},p^{(k)})},\\
w^{(k)}&=w^{(k-1)}+\alpha_kp^{(k)},\\
r^{(k)}&=r^{(k-1)}-\alpha_kAp^{(k)},\qquad
z^{(k)}=D_{\rm diag}^{-1}r^{(k)},\\
\beta_{k+1}&=\frac{(r^{(k)},z^{(k)})}{(r^{(k-1)},z^{(k-1)})},\qquad
p^{(k+1)}=z^{(k)}+\beta_{k+1}p^{(k)}.
\end{align*}
Критерий остановки:
$\norm{w^{(k)}-w^{(k-1)}}_E<\delta$, где $\delta=10^{-9}$.
Перед принятием результата заново вычисляется истинная невязка; дополнительно требуется
$\norm{B-Aw}_E/\norm B_E<10^{-8}$. Число итераций ограничено числом неизвестных.

Контролируется функционал
\[
J(w)=\tfrac12(Aw,w)-(B,w)=-\tfrac12(B+r,w).
\]
При вычислении суммы для $J$ применяется компенсация Кэхэна внутри каждой нити. Рост выше допуска округления приводит к откату шага и перезапуску с истинной невязкой. В окончательной серии перезапуски не потребовались.

Исходный код написан на C++11. Сеточные функции хранятся в векторах \texttt{std::vector}, которые автоматически освобождают память. Узел $(i,j)$ имеет индекс $i(N+1)+j$. Время измеряется монотонными часами \texttt{steady\_clock} из библиотеки \texttt{chrono}. Файлы записываются через потоки C++. Неверные аргументы и ошибки чтения файлов приводят к диагностике и ненулевому коду возврата.

\section{Последовательная версия и проверка точности}
Расчёты выполнены на сетках $20\times20$, $40\times40$, $80\times80$. В таблице приведены медианы пяти измерений после прогрева.
\begin{table}[H]
\centering\small
\caption{Последовательная C++-версия, медианы пяти измерений.}
\begin{tabular}{rrrrrr}
\toprule
$M\times N$ & Неизвестных & $\varepsilon$ & Итераций & Время, с & $\|r\|_E/\|B\|_E$ \\
\midrule
$20\times20$ & 361 & $0{,}09$ & 70 & $0{,}000923668$ & $7{,}44\cdot10^{-11}$ \\
$40\times40$ & 1521 & $0{,}0225$ & 134 & $0{,}00735755$ & $2{,}12\cdot10^{-11}$ \\
$80\times80$ & 6241 & $0{,}005625$ & 261 & $0{,}0581861$ & $4{,}26\cdot10^{-12}$ \\
\bottomrule
\end{tabular}
\end{table}


\subsection{Тест с известным точным решением}
В отдельном тесте правая часть заменена на $F(x,y)=2\sin(x+y)$ во всём прямоугольнике, включая фиктивную область. При прежнем $v_0$ точным непрерывным решением служит $u_{\rm exact}=\sin(x+y)$. Для теста использовано $\delta=10^{-11}$.
\begin{table}[H]
\centering\small
\caption{Погрешность в проверочной задаче с точным решением.}
\begin{tabular}{rrrr}
\toprule
$M\times N$ & $\|w-u_{\rm exact}\|_E$ & $\|w-u_{\rm exact}\|_\infty$ & Отношение норм \\
\midrule
$20\times20$ & $0{,}01476$ & $0{,}00652957$ & --- \\
$40\times40$ & $0{,}00352337$ & $0{,}00160082$ & $4{,}189$ \\
$80\times80$ & $0{,}000873062$ & $0{,}000399679$ & $4{,}036$ \\
\bottomrule
\end{tabular}
\end{table}

При сгущении сетки вдвое ошибка уменьшается примерно в четыре раза, что согласуется со вторым порядком разностной аппроксимации. Эти ошибки относятся к проверочной задаче, а не к исходной задаче с $f=1$.

\subsection{Тесты программы}
Команда \texttt{make test} собирает и запускает C++-тесты из \texttt{tests/check.cpp}. В тестах проверяются:
\begin{itemize}
\setlength{\itemsep}{0pt}
\item площадь всей области по независимой геометрической формуле
$|D|=9\pi-18\arccos(2/3)+4\sqrt5$ с допуском $10^{-11}$;
\item площади целиком внутренней и целиком внешней ячеек;
\item границы $0\leq S\leq |\Pi_{ij}|$ и аддитивность площади для 300 воспроизводимых случайных ячеек;
\item убывание ошибки в задаче с точным решением на трёх сетках;
\item совпадение решений на 1, 4 и 16 нитях с допуском $10^{-10}$;
\item по одной записи истории на итерацию, выполнение критерия остановки и отсутствие существенного роста функционала.
\end{itemize}
Тесты выполнены на Polus тем же компилятором с ключом \texttt{-O1}. Результаты проверки сохранены в \texttt{results/tests.log}.
\begin{table}[H]
\centering\small
\caption{Совпадение последовательной и OpenMP-версий на сетке $80\times80$.}
\begin{tabular}{rr}
\toprule
Нити & $\max|w_{\rm OpenMP}-w_{\rm serial}|$ \\
\midrule
1 & $0$ \\
4 & $6\cdot10^{-15}$ \\
16 & $5{,}33\cdot10^{-15}$ \\
\bottomrule
\end{tabular}
\end{table}


\section{Сходимость и ошибка по итерациям}
В диагностических запусках на каждой итерации заново вычисляется истинная невязка $r_{\rm true}^{(k)}=B-Aw^{(k)}$. Это позволяет проверять рекуррентное обновление невязки независимо от накопления округлений. На рисунке показаны норма изменения решения и относительная истинная невязка. Небольшая невязка относится к разностной системе и сама по себе не является оценкой ошибки исходной краевой задачи.
\pic{convergence.png}{1.0}{Изменение решения и истинная невязка по итерациям для исходной задачи.}

\subsection{Итерационная ошибка исходной задачи}
Для исходной задачи с $f=1$ точное решение неизвестно. Поэтому на каждой из тех же сеток отдельно вычислено решение $w_{\rm ref}$ с более строгим допуском $\delta_{\rm ref}=10^{-12}$. На рисунке показаны
\[
e_E^{(k)}=\norm{w^{(k)}-w_{\rm ref}}_E,\qquad
e_\infty^{(k)}=\max_{1\leq i<M,\,1\leq j<N}|w_{ij}^{(k)}-(w_{\rm ref})_{ij}|.
\]
Это оценка итерационной ошибки относительно более точного решения той же разностной системы. Она не включает ошибку сеточной аппроксимации и метода фиктивных областей; само $w_{\rm ref}$ также вычислено приближённо.
\pic{iteration_error.png}{1.0}{Итерационная ошибка исходной задачи в нормах $E$ и $\infty$.}
\begin{table}[H]
\centering\small
\caption{Итерационная ошибка при остановке относительно $w_{\rm ref}$ с $\delta=10^{-12}$.}
\begin{tabular}{rrr}
\toprule
$M\times N$ & $\|w-w_{\rm ref}\|_E$ & $\|w-w_{\rm ref}\|_\infty$ \\
\midrule
$20\times20$ & $7{,}608\cdot10^{-10}$ & $4{,}617\cdot10^{-10}$ \\
$40\times40$ & $2{,}822\cdot10^{-9}$ & $1{,}835\cdot10^{-9}$ \\
$80\times80$ & $5{,}541\cdot10^{-9}$ & $3{,}346\cdot10^{-9}$ \\
\bottomrule
\end{tabular}
\end{table}


\subsection{Ошибка относительно известного точного решения}
Для проверочной задачи с $u_{\rm exact}=\sin(x+y)$ строится отдельный график $\max|w^{(k)}-u_{\rm exact}|$. После снижения итерационной ошибки график выходит на уровень погрешности разностной схемы. Поэтому уменьшение допуска итерационного решателя не может устранить всю ошибку непрерывной задачи на фиксированной сетке.
\pic{exact_error.png}{0.96}{Ошибка по итерациям в проверочной задаче с известным решением.}

\subsection{Сходимость в разных частях области}
На сетке $80\times80$ сохранены промежуточные решения через каждые 50 итераций и в момент остановки. Сначала показано формирование самого решения с общей цветовой шкалой.
\pic{solution_iterations.png}{1.0}{Приближения $w^{(k)}$ на разных итерациях, сетка $80\times80$.}

Для каждого снимка также построено поле
$|w_{ij}^{(k)}-(w_{\rm ref})_{ij}|$ внутри $D$. Все четыре изображения используют одну логарифмическую цветовую шкалу; значения ниже $10^{-10}$ показаны нижним цветом шкалы. Эти карты позволяют сопоставлять скорость уменьшения ошибки в разных участках области.
\pic{spatial_error.png}{1.0}{Пространственное распределение итерационной ошибки на разных итерациях. Фиктивная область скрыта.}

Дополнительно для каждого узла найден первый сохранённый снимок, начиная с которого ошибка на всех оставшихся снимках не превышает $10^{-5}$. Тёмные участки достигают этого уровня раньше, светлые --- позже. Номера получены только по сохранённым снимкам: это оценка с интервалом 50 итераций, а не точный номер первого пересечения порога. Между снимками локальная ошибка может вести себя немонотонно.
\pic{spatial_convergence.png}{0.83}{Неоднородность сходимости: достижение порога $10^{-5}$ по сохранённым снимкам.}

\section{Визуализация решения}
Фиктивная область на рисунках скрыта; линия показывает границу $D$.
\pic{solution_800_1200.png}{0.75}{Приближённое решение на самой большой сетке $800\times1200$.}
Постоянная правая часть не делает решение симметричным относительно оси $y$: граничная функция $\sin(x+y)$ такой симметрией не обладает.

\section{OpenMP-реализация}
В файле \texttt{src/solver.cpp} последовательно выполняются итерации метода, а внутри каждой итерации распараллелены независимые вычисления по внутренним узлам: построение коэффициентов, применение оператора и обновление векторов. Использованы директивы \texttt{parallel for}, \texttt{collapse(2)}, \texttt{schedule(static)} и редукции для скалярных произведений и норм. Каждая нить записывает свои элементы массивов.

\subsection{Основная сетка}
\begin{table}[H]
\centering\small
\caption{OpenMP на основной сетке.}
\begin{tabular}{rrrrr}
\toprule
Нити & $M\times N$ & Итераций & Время, с & $S$ \\
\midrule
1 & $80\times80$ & 261 & $0{,}0776358$ & $0{,}7495$ \\
4 & $80\times80$ & 261 & $0{,}0257036$ & $2{,}264$ \\
16 & $80\times80$ & 261 & $0{,}0257699$ & $2{,}258$ \\
\bottomrule
\end{tabular}
\end{table}

\subsection{Таблица 1 из приложения}
\begin{table}[H]
\centering\small
\caption{Последовательная база больших сеток в том же задании.}
\begin{tabular}{rrr}
\toprule
$M\times N$ & Итераций & Время, с \\
\midrule
$400\times600$ & 1638 & $13{,}9847$ \\
$800\times1200$ & 3178 & $109{,}283$ \\
\bottomrule
\end{tabular}
\end{table}

\begin{table}[H]
\centering\small
\caption{Таблица 1 приложения: OpenMP, ключ -O1.}
\begin{tabular}{rrrrr}
\toprule
Нити & $M\times N$ & Итераций & Время, с & $S$ \\
\midrule
2 & $400\times600$ & 1638 & $9{,}37306$ & $1{,}492$ \\
4 & $400\times600$ & 1638 & $4{,}73745$ & $2{,}952$ \\
8 & $400\times600$ & 1638 & $2{,}42671$ & $5{,}763$ \\
16 & $400\times600$ & 1638 & $1{,}38297$ & $10{,}11$ \\
4 & $800\times1200$ & 3178 & $36{,}6$ & $2{,}986$ \\
8 & $800\times1200$ & 3178 & $18{,}3922$ & $5{,}942$ \\
16 & $800\times1200$ & 3178 & $9{,}47854$ & $11{,}53$ \\
32 & $800\times1200$ & 3178 & $7{,}23031$ & $15{,}11$ \\
\bottomrule
\end{tabular}
\end{table}

\pic{speedup_openmp.png}{1.0}{Измеренное ускорение OpenMP относительно последовательной версии на той же сетке.}
На малой сетке накладные расходы параллельных областей и синхронизаций ограничивают ускорение. На больших сетках полезной работы на нить больше. На выделенном узле имеется 20 физических ядер. Для 16 нитей использовано 16 различных ядер, равномерно выбранных из двух сокетов. Для 32 нитей на тех же 16 ядрах использовано по два аппаратных потока (SMT). Ускорение для 32 нитей делится на 32 при вычислении эффективности; 32 нити не означают 32 физических ядра.

\section{Условия измерений и воспроизводимость}
Расчёты выполнены 9 ноября 2026 года для C++-версии с ключом \texttt{-O1}.

Измерения выполнены на одном эксклюзивно выделенном узле IBM Polus, \texttt{polus-c2}, в одном задании LSF 1614435. Процессоры POWER8NVL, архитектура \texttt{ppc64le}; два сокета по 10 физических ядер, SMT8. Компилятор G++ 4.8.5, параметры
\texttt{-O1 -std=c++11 -Wall -Wextra}; для параллельной версии добавлен \texttt{-fopenmp}.

Нити закреплены через \texttt{GOMP\_CPU\_AFFINITY}; динамическое изменение размера команды отключено через \texttt{OMP\_DYNAMIC=FALSE}. Использованы первые аппаратные потоки выбранных ядер, а для 32 нитей --- ещё по одному потоку тех же ядер. Последовательная программа закреплена за одним из этих ядер командой \texttt{taskset}. Точный список CPU, сведения об узле и компиляторе сохранены вместе с результатами.

На малых сетках выполнено пять измерений на конфигурацию, на больших --- три; таблицы содержат медианы. Перед каждой конфигурацией выполнен отдельный прогрев. Всего для оценки времени выполнено 60 измерений в 16 конфигурациях. Последовательная база для каждого ускорения взята с той же сетки, того же узла и из того же задания:
\[
S_p=\frac{T_{\rm serial}}{T_{{\rm OpenMP},p}},\qquad E_p=\frac{S_p}{p}.
\]

Таймер охватывает итерационный решатель, включая пересчёт истинной невязки при остановке. Он исключает очередь, построение коэффициентов, чтение эталона и итоговую запись решения. История и снимки получены в отдельных диагностических запусках; во всех замерах времени они отключены. Если история или снимки включены вручную, их вычисление и запись внутри итераций входят в печатаемое время, поэтому такие времена не используются для ускорения.

\subsection{Сборка и запуск}
Из каталога лабораторной команда \texttt{make} собирает \texttt{serial} и \texttt{openmp}, а \texttt{make test} запускает тесты. Флаг \texttt{-O1} задан по умолчанию в \texttt{Makefile}. Файл \texttt{bench.lsf} содержит воспроизводимую серию и подаётся на вычислительный узел командой \texttt{bsub < bench.lsf}.

Пример запуска параллельной программы на уже выделенном вычислительном узле:
\begin{quote}\ttfamily\small
OMP\_NUM\_THREADS=4 ./openmp 80 80 1e-9 - -
\end{quote}
Аргументы: $M$, $N$, допуск, файл итогового решения (или \texttt{-}), файл истории (или \texttt{-}), необязательный флаг проверочной задачи 0/1. Ключ \texttt{--reference} задаёт файл более точного решения той же сетки, \texttt{--snapshots} --- существующий каталог снимков, \texttt{--snapshot-every} --- интервал записи. Для перестроения таблиц и рисунков используется \texttt{python3 analyze.py}; требуются NumPy и Matplotlib. Большое итоговое решение допускается хранить в \texttt{CSV.gz}; скрипт читает оба формата.

\section{Выводы}
В первом этапе разработаны две C++-версии программы: последовательная и с OpenMP. Обе используют одну разностную схему. В тесте с точным решением удвоение числа интервалов уменьшает ошибку примерно в четыре раза.

Решения обеих версий совпадают в пределах округления. Во всех окончательных расчётах первого этапа достигнут критерий остановки. Графики показывают различие между невязкой, итерационной ошибкой и ошибкой аппроксимации; пространственные карты позволяют сравнить сходимость в разных частях области.

Наибольшее ускорение в измеренной серии составило $15{,}11$ при 32 нитях на сетке $800\times1200$.


\begingroup
\small
\begin{thebibliography}{9}
\bibitem{task} Мокин А.\,Ю. Задание по курсу «Суперкомпьютерное моделирование и технологии». Сентябрь 2026. Разделы 3--6 и приложение.
\bibitem{variants} Распределение вариантов.xlsx. Строка: Кабедев Вадим Дмитриевич, группа 609, задача 1, вариант 15.
\bibitem{polus} Высокопроизводительные вычисления на ВМК МГУ. Подключение и запуск заданий Polus: \url{https://hpc.cs.msu.ru/node/245}, \url{https://hpc.cs.msu.ru/node/242}.
\end{thebibliography}
\endgroup
\end{document}
